% !TEX program = pdflatex
% Overleaf-ready single-file report.
% Compile with pdflatex on Overleaf or locally.

\documentclass[11pt]{article}

\usepackage[margin=1in]{geometry}
\usepackage{hyperref}
\usepackage{listings}
\usepackage{xcolor}
\usepackage{enumitem}
\usepackage{parskip}
\usepackage{graphicx}
\usepackage{tabularx}
\usepackage{booktabs}
\usepackage{titlesec}

\hypersetup{
  colorlinks=true,
  linkcolor=black,
  urlcolor=blue,
  citecolor=black
}

\definecolor{codebg}{RGB}{246,246,250}
\definecolor{codekw}{RGB}{170,40,90}
\definecolor{codecmt}{RGB}{100,110,120}
\definecolor{codestr}{RGB}{30,120,80}

\lstdefinelanguage{JavaScriptPlus}{
  keywords={break, case, catch, const, continue, debugger, default, delete, do,
    else, false, finally, for, function, if, in, instanceof, let, new, null,
    return, switch, this, throw, true, try, typeof, var, void, while, with,
    yield, async, await, of, export, import, from, class, extends, super,
    static, true, false},
  keywordstyle=\color{codekw}\bfseries,
  ndkeywords={describe, it, expect, beforeEach, afterEach, jasmine,
    Given, When, Then, async},
  ndkeywordstyle=\color{codekw}\bfseries,
  identifierstyle=\color{black},
  sensitive=true,
  comment=[l]{//},
  morecomment=[s]{/*}{*/},
  commentstyle=\color{codecmt}\itshape,
  stringstyle=\color{codestr},
  morestring=[b]',
  morestring=[b]"
}

\lstset{
  language=JavaScriptPlus,
  basicstyle=\ttfamily\footnotesize,
  backgroundcolor=\color{codebg},
  frame=single,
  framerule=0pt,
  rulecolor=\color{codebg},
  breaklines=true,
  showstringspaces=false,
  columns=fullflexible,
  xleftmargin=6pt,
  xrightmargin=6pt,
  aboveskip=8pt,
  belowskip=8pt
}

\titleformat{\section}{\large\bfseries}{\thesection}{0.6em}{}
\titleformat{\subsection}{\normalsize\bfseries}{\thesubsection}{0.5em}{}

\title{Individual Iteration 1\\ \large Multi-LLM Local Compare on Scarlet AI}
\author{Partha Pediredla \\ NetID: pp955 \\ GitHub: \texttt{parthaped} \\ Software Engineering, Spring 2026, Group 1}
\date{April 23, 2026}

\begin{document}
\maketitle

\section*{Submission Pointers}

\begin{tabularx}{\textwidth}{lX}
\toprule
Repository & \url{https://github.com/Rutgers-SP26-SWE-Group-1/SP26SWE---Group1} \\
Branch & \texttt{parthaped-individual-report} \\
Branch URL & \url{https://github.com/Rutgers-SP26-SWE-Group-1/SP26SWE---Group1/tree/parthaped-individual-report} \\
Local feature page & \url{http://localhost:3000/chat/multi} \\
Customer need addressed & ``Get answers from multiple LLMs simultaneously.'' \\
LLM constraint (mine) & Local Ollama models only (no cloud APIs in any new code path). \\
\bottomrule
\end{tabularx}

\section{User Stories and Features}

\subsection{User Stories}

\paragraph{User Story 1 (3 points).}
As a Rutgers student, so that I can compare how different small open models
phrase the same answer, I want to type one prompt on a dedicated compare page
and have every model I selected reply at the same time.

\paragraph{User Story 2 (2 points).}
As a Rutgers student running models on my own laptop, so that I do not waste
disk or RAM on models I am not using right now, I want to pick which local
models answer my next question from a chip-style dropdown that caps the
selection at four.

\paragraph{User Story 3 (2 points).}
As a Rutgers student, so that I can quickly tell which model gave the
fastest or clearest answer, I want every reply rendered in its own card
inline on the same page, with a model label, a latency badge, and an obvious
``fastest'' marker on the quickest successful response.

\subsection{SMART Justification}

\begin{itemize}[leftmargin=*]
  \item \textbf{Specific.} Each story names a single screen
    (\texttt{/chat/multi}), a single UI primitive (chip dropdown,
    inline grid), and a single backend route (\texttt{POST
    /api/chat/fanout}).
  \item \textbf{Measurable.} The selection cap is a literal constant
    (\texttt{MAX\_MODELS = 4}); the per-card latency badge is asserted in
    Puppeteer as a regex \texttt{/\textbackslash d+\textbackslash.\textbackslash d+s/}.
  \item \textbf{Achievable.} The work reuses the team's existing
    \texttt{validateChatRequest}, \texttt{sanitizeMessages}, and
    \texttt{createConversationId} helpers; no new third-party SDK is
    introduced.
  \item \textbf{Relevant.} Directly answers the customer need without
    duplicating teammate Shiv's modal-based ``select 3 models'' work.
  \item \textbf{Timeboxed.} Scope sized for a single iteration; only the
    new files under \texttt{src/lib/multi-llm/}, \texttt{src/components/multi/},
    \texttt{src/app/api/chat/fanout/}, and \texttt{src/app/chat/multi/}
    are touched.
\end{itemize}

\subsection{Selected Features}

\begin{enumerate}[leftmargin=*]
  \item \textbf{Feature 1: Local-models chip dropdown} with a 1--4 cap and
    a live ``\(n/4\)'' counter, backed by a typed local registry that
    knows nothing about cloud providers.
  \item \textbf{Feature 2: Parallel server-side fan-out} of one prompt to
    every selected local model via a dedicated
    \texttt{POST /api/chat/fanout} route that uses \texttt{Promise.all}
    with per-model error isolation.
  \item \textbf{Feature 3: Inline side-by-side comparison grid} that
    appears in the same page as the prompt composer (no overlay, no
    modal), with a fastest-response badge and a copy-to-clipboard
    button per card.
\end{enumerate}

\subsection{What Makes This Submission Individualized}

\begin{itemize}[leftmargin=*]
  \item \textbf{Different UI primitive.} Chip dropdown anchored beside the
    composer instead of a floating modal card.
  \item \textbf{Different backend shape.} A new \texttt{/api/chat/fanout}
    route that runs all calls server-side in parallel, returning a single
    \texttt{\{ responses: [...] \}} payload. Other team members loop on
    the client.
  \item \textbf{Different model set.} Five local Ollama models with tags
    that are deliberately not the same as teammate Shiv's
    \texttt{mistral / llama3.1 / deepseek-r1:8b / qwen2.5-coder:7b /
    gemma3:4b}.  The set used here is
    \texttt{llama3.2:latest}, \texttt{phi3:mini},
    \texttt{qwen2.5:3b}, \texttt{tinyllama:latest},
    \texttt{gemma2:2b}.
  \item \textbf{Different render path.} An inline CSS-grid of cards on
    the same page, not a separate modal that opens and closes.
  \item \textbf{Local-only.} The new code path imports nothing from
    teammates' cloud-routed handler, so the feature cannot accidentally
    fall back to Gemini, Groq, or Anthropic.
\end{itemize}

\section{UI Design}

\subsection{Pages in the Project}

\begin{tabularx}{\textwidth}{lXl}
\toprule
\textbf{Route} & \textbf{Purpose} & \textbf{Owner} \\
\midrule
\texttt{/} & Marketing landing & shared \\
\texttt{/login}, \texttt{/signup} & Supabase auth (\texttt{@scarletmail.rutgers.edu}) & shared \\
\texttt{/profile} & User profile management & shared \\
\texttt{/chat} & Single-model chat hub & shared \\
\texttt{/chat/multi} & Local multi-LLM compare (\textbf{new in this iteration}) & Partha \\
\bottomrule
\end{tabularx}

\subsection{User Flow}

\begin{enumerate}[leftmargin=*]
  \item User signs in (or skips auth) and lands on \texttt{/chat}.
  \item User clicks ``Local multi-LLM compare'' in the chat hub
    (or visits \texttt{/chat/multi} directly).
  \item User opens the chip-style dropdown and confirms a 1--4 model
    selection.  The selection is persisted in
    \texttt{localStorage} under
    \texttt{scarlet-ai-multi-llm-selection}.
  \item User types a single prompt in the composer and presses
    ``Ask all selected''.
  \item The browser POSTs to \texttt{/api/chat/fanout}; the server
    fans out to every selected Ollama model in parallel and returns
    one combined payload.
  \item The page renders a card per model (label, family, Ollama
    tag, latency, copy button) inline in the transcript.  The fastest
    successful card receives a green ``fastest'' badge.
  \item User may clear the transcript or repeat with a new prompt or a
    new selection.
\end{enumerate}

\subsection{Lo-Fi Sketches}

\begin{quote}\itshape
The lo-fi sketches done in class are stored in
\texttt{report/figures/lofi-feature\{1,2,3\}.png}.  If the figures are
not present at compile time the placeholders below render as gray
boxes; replace the \texttt{\textbackslash includegraphics} calls
once the scans are added.
\end{quote}

\begin{center}
\fbox{\begin{minipage}{0.92\textwidth}
\centering\vspace{0.6cm}
\textit{Figure placeholder:} \texttt{report/figures/lofi-feature1.png}\\
Chip dropdown anchored next to the prompt composer; cap badge shows ``3/4''.
\vspace{0.6cm}
\end{minipage}}
\end{center}

\begin{center}
\fbox{\begin{minipage}{0.92\textwidth}
\centering\vspace{0.6cm}
\textit{Figure placeholder:} \texttt{report/figures/lofi-feature2.png}\\
One prompt $\rightarrow$ fan-out to all selected models $\rightarrow$ one merged payload.
\vspace{0.6cm}
\end{minipage}}
\end{center}

\begin{center}
\fbox{\begin{minipage}{0.92\textwidth}
\centering\vspace{0.6cm}
\textit{Figure placeholder:} \texttt{report/figures/lofi-feature3.png}\\
Inline grid of model cards on the same page; ``fastest'' badge on the quickest reply.
\vspace{0.6cm}
\end{minipage}}
\end{center}

\section{Unit Tests and Acceptance Tests}

\subsection{TDD Procedure}

For each feature I wrote the Jasmine spec first against a pure helper
module, then implemented the module, then wired the module into the
React component or Next.js route, and finally added the Cucumber and
Puppeteer flows.  Concretely:

\begin{enumerate}[leftmargin=*]
  \item Draft the user story.
  \item Write Jasmine \texttt{describe}/\texttt{it} blocks describing
    the expected pure-function behaviour.
  \item Implement the helper in \texttt{src/lib/multi-llm/} until the
    spec passes.
  \item Write Cucumber scenarios that drive the same behaviour through
    the live UI.
  \item Implement the React component / API route until both layers
    pass.
  \item Add a Puppeteer script as the headed acceptance demo.
  \item Commit and push.
\end{enumerate}

\subsection{Test-Suite Layout}

\begin{itemize}[leftmargin=*]
  \item \textbf{Jasmine} (\texttt{spec/}): one spec file per feature,
    each one a single \texttt{describe} block grouping all
    \texttt{it} cases for that feature.  Pure-function tests; no
    network, no DOM.
  \item \textbf{Cucumber} (\texttt{features/}): one feature file per
    feature.  Tagged with \texttt{@multi1}, \texttt{@multi2},
    \texttt{@multi3} so each can be invoked independently.  Steps
    live in
    \texttt{features/step\_definitions/multiSteps.js}.
  \item \textbf{Puppeteer} (\texttt{tests/}): one script per feature.
    Each opens a browser, seeds the localStorage selection, and asserts
    DOM and network behaviour, dumping screenshots into
    \texttt{test-screenshots/}.
\end{itemize}

\section{Feature 1: Local-Models Chip Dropdown}

\subsection{Cucumber Acceptance Scenarios}

\begin{lstlisting}
@multi @multi1
Feature: Multi-LLM dropdown selection (Feature 1, parthaped)

  Background:
    Given the user is on the multi-LLM compare page

  Scenario: The chip-style dropdown opens on click
  Scenario: The user can add a fourth distinct local model
  Scenario: The dropdown refuses to add a fifth local model
  Scenario: The dropdown refuses to drop below the one-model floor
\end{lstlisting}

Full file: \texttt{features/multi-select.feature}.  Steps:
\texttt{features/step\_definitions/multiSteps.js}.

\subsection{Jasmine Unit Tests}

Spec file: \texttt{spec/multi\_selection\_spec.js}.  All nine cases
exercise the pure module
\texttt{src/lib/multi-llm/selection.js}:

\begin{itemize}[leftmargin=*]
  \item exposes the \(1\le n \le 4\) selection cap;
  \item seeds the first registered local model when no initial ids
    are supplied;
  \item honours a caller-supplied seed but caps it at
    \texttt{MAX\_MODELS} and ignores unknown ids;
  \item opens and closes the dropdown via pure transitions;
  \item lets the user add a second, third, and fourth distinct model;
  \item refuses to add a fifth model once the cap is reached;
  \item lets the user deselect down to one model but no further;
  \item silently ignores toggling an unknown model id;
  \item \texttt{clearSelection} resets to a single default model.
\end{itemize}

Representative snippet:

\begin{lstlisting}
it('refuses to add a fifth model once the cap is reached', function () {
  let state = createSelectionState(LOCAL_OLLAMA_MODELS.slice(0, 4).map((m) => m.id));
  const before = state.selectedModelIds.slice();
  state = toggleModel(state, LOCAL_OLLAMA_MODELS[4].id);
  expect(state.selectedModelIds).toEqual(before);
  expect(selectionCount(state)).toBe(MAX_MODELS);
});
\end{lstlisting}

\subsection{Puppeteer}

Script: \texttt{tests/parthaped-feature1-puppeteer-test.js}.
Asserts: page load; default ``3/4'' counter; dropdown opens on click;
all five local models listed; fourth selectable; fifth disabled at
the cap; deselection floor at 1; the lone selected option becomes
unclickable.

\section{Feature 2: Parallel Fan-Out}

\subsection{Cucumber Acceptance Scenarios}

\begin{lstlisting}
@multi @multi2
Feature: Parallel fan-out to selected local models

  Scenario: One prompt produces one card per selected model
  Scenario: Each card is labelled with the local model that produced it
  Scenario: A failure in one model does not block the others
\end{lstlisting}

Full file: \texttt{features/multi-fanout.feature}.

\subsection{Jasmine Unit Tests}

Spec file: \texttt{spec/multi\_fanout\_spec.js}.  Coverage:

\begin{itemize}[leftmargin=*]
  \item rejects an empty model list with a useful error;
  \item rejects an empty prompt message list;
  \item dispatches the same prompt to every selected model;
  \item produces one response per selected model with model id and
    label;
  \item isolates per-model failures so siblings still resolve;
  \item runs requests concurrently rather than sequentially
    (timed assertion);
  \item records a non-negative \texttt{durationMs} for every
    response;
  \item \texttt{summarizeFanout} reports correct succeeded/failed
    counts;
  \item returns an error entry (not a throw) when an unknown model id
    sneaks in.
\end{itemize}

The concurrency check is the most important assertion of the feature:

\begin{lstlisting}
it('runs requests concurrently rather than sequentially', async function () {
  const startedAt = Date.now();
  const PER_MODEL_DELAY_MS = 80;
  await runFanout({
    modelIds: [
      LOCAL_OLLAMA_MODELS[0].id,
      LOCAL_OLLAMA_MODELS[1].id,
      LOCAL_OLLAMA_MODELS[2].id,
      LOCAL_OLLAMA_MODELS[3].id,
    ],
    promptMessages: PROMPT_MESSAGES,
    requestModelResponse: ({ option }) =>
      new Promise((resolve) => {
        setTimeout(() => resolve(`${option.label} done`), PER_MODEL_DELAY_MS);
      }),
  });
  const elapsed = Date.now() - startedAt;
  expect(elapsed).toBeLessThan(PER_MODEL_DELAY_MS * 3);
});
\end{lstlisting}

\subsection{Puppeteer}

Script: \texttt{tests/parthaped-feature2-puppeteer-test.js}.
Asserts: seeded 2/4 selection; one card per selected model after
submission; every card has a known local-model label; every card
shows a latency reading; injecting an invalid model id produces
exactly one error card and the valid models still respond.

\section{Feature 3: Inline Side-by-Side Grid}

\subsection{Cucumber Acceptance Scenarios}

\begin{lstlisting}
@multi @multi3
Feature: Side-by-side comparison grid

  Scenario: The grid renders one column per model on smaller selections
  Scenario: The grid stays visible after the answer arrives
  Scenario: Clearing the transcript empties the grid
\end{lstlisting}

Full file: \texttt{features/multi-grid.feature}.

\subsection{Jasmine Unit Tests}

Spec file: \texttt{spec/multi\_response\_grid\_spec.js}.  Coverage:

\begin{itemize}[leftmargin=*]
  \item returns an empty array when there are no responses;
  \item produces one section per response, in submission order;
  \item resolves human-friendly labels and family names from the
    registry;
  \item falls back to the response label when the model id is unknown;
  \item marks rejected responses as errors and surfaces the message;
  \item \texttt{pickFastestSection} picks the lowest \texttt{durationMs}
    among successful responses;
  \item \texttt{pickFastestSection} ignores rejected responses;
  \item \texttt{pickFastestSection} returns null when every response
    failed;
  \item \texttt{gridColumnCount} ramps \(1\to1, 2\to2, 3+\to3\)
    columns.
\end{itemize}

Representative snippet:

\begin{lstlisting}
it('pickFastestSection ignores rejected responses', function () {
  const sections = buildGridSections([
    makeRejected(LOCAL_OLLAMA_MODELS[0].id, 'oops', 5),
    makeFulfilled(LOCAL_OLLAMA_MODELS[1].id, 'b', 700),
  ]);
  const fastest = pickFastestSection(sections);
  expect(fastest.modelId).toBe(LOCAL_OLLAMA_MODELS[1].id);
});
\end{lstlisting}

\subsection{Puppeteer}

Script: \texttt{tests/parthaped-feature3-puppeteer-test.js}.
Asserts: the grid is hidden before any prompt is sent; appears inline
after submission; one card per selected model; every card has a
distinct model label and a latency reading; exactly one card is
marked ``fastest''; the grid is removed after clearing the
transcript.

\section{Software Architecture and Implementation}

\subsection{REST API}

The new code adds exactly one endpoint:

\begin{lstlisting}
POST /api/chat/fanout

Request body
{
  "message":        "<user prompt>",          // required, <= 2000 chars
  "modelIds":       ["llama3.2", "phi3-mini"],// 1..4 entries, all local
  "messages":       [ { role, content }, ... ], // optional history
  "conversationId": "<existing-uuid>"         // optional
}

Response (200)
{
  "conversationId": "<uuid>",
  "responses": [
    {
      "modelId": "llama3.2",
      "modelLabel": "Llama 3.2",
      "ollamaTag": "llama3.2:latest",
      "status": "fulfilled" | "rejected",
      "content": "...",
      "error":   "..." (only when rejected),
      "durationMs": 832
    },
    ...
  ],
  "summary":          { "total":3, "succeeded":2, "failed":1, "averageMs":910 },
  "totalDurationMs":  1240,
  "timestamp":        "2026-04-23T22:00:00.000Z"
}

Error (400) when message empty / missing / >2000 chars or no valid local
model ids; (500) when the fan-out itself throws.
\end{lstlisting}

\subsection{Routing Table}

\begin{tabularx}{\textwidth}{lllX}
\toprule
\textbf{Method} & \textbf{Path} & \textbf{Status} & \textbf{Purpose} \\
\midrule
GET  & \texttt{/} & shared & Landing page \\
GET  & \texttt{/login}, \texttt{/signup} & shared & Supabase auth screens \\
POST & \texttt{/api/auth/signup} & shared & Sign-up handler \\
GET  & \texttt{/profile} & shared & Profile management \\
GET  & \texttt{/chat} & shared & Single-model chat hub \\
GET  & \texttt{/chat/multi} & \textbf{new} & Multi-LLM compare page (Partha) \\
POST & \texttt{/api/chat} & shared & Single-model chat (existing) \\
POST & \texttt{/api/chat/fanout} & \textbf{new} & Parallel local fan-out (Partha) \\
\bottomrule
\end{tabularx}

\subsection{Database Design}

The shared Supabase schema is unchanged: it stores a row per
authenticated user (\texttt{@scarletmail.rutgers.edu}) plus profile
metadata.  Multi-model conversations themselves are deliberately
\emph{not} persisted server-side this iteration; they live in
\texttt{localStorage}:

\begin{tabularx}{\textwidth}{lX}
\toprule
\textbf{Key} & \textbf{Shape} \\
\midrule
\texttt{scarlet-ai-multi-llm-selection} &
\texttt{string[]} of local-model ids that survives reload. \\
\texttt{scarlet-ai-conversations} (existing) &
Single-model transcripts kept by teammates' chat hub. \\
\bottomrule
\end{tabularx}

The fan-out route reuses
\texttt{createConversationId} from \texttt{src/lib/chat-logic.js},
so when a future iteration starts persisting multi-model turns into
Supabase the response payload already carries a stable id.

\subsection{Component / Module Layout}

\begin{itemize}[leftmargin=*]
  \item \texttt{src/lib/multi-llm/localModels.js} -- registry of the
    five local Ollama models (id, label, ollamaTag, family, blurb,
    sizeMb), plus \texttt{getLocalModel}, \texttt{listLocalModelIds},
    \texttt{isLocalModelId}.
  \item \texttt{src/lib/multi-llm/selection.js} -- pure dropdown state
    machine.
  \item \texttt{src/lib/multi-llm/ollamaClient.js} -- thin
    \texttt{POST /api/chat} client; throws typed
    \texttt{OllamaUnavailableError} or
    \texttt{OllamaModelMissingError}.
  \item \texttt{src/lib/multi-llm/fanout.js} --
    \texttt{Promise.all}-based fan-out with per-model error isolation;
    \texttt{requestModelResponse} is dependency-injected for testability.
  \item \texttt{src/lib/multi-llm/grid.js} --
    \texttt{buildGridSections}, \texttt{pickFastestSection},
    \texttt{gridColumnCount}.
  \item \texttt{src/app/api/chat/fanout/route.ts} -- Next.js route
    handler; imports nothing from the cloud-routed
    \texttt{src/app/api/chat/route.ts}.
  \item \texttt{src/components/multi/MultiModelDropdown.tsx} --
    chip dropdown.
  \item \texttt{src/components/multi/MultiResponseGrid.tsx} +
    \texttt{ModelResponseCard.tsx} -- inline CSS-grid renderer.
  \item \texttt{src/app/chat/multi/page.tsx} -- the new feature page.
\end{itemize}

\section{Instructions for the TAs}

\subsection{Installation}

\begin{lstlisting}
git clone https://github.com/Rutgers-SP26-SWE-Group-1/SP26SWE---Group1.git
cd SP26SWE---Group1
git fetch origin
git checkout -b parthaped-individual-report origin/parthaped-individual-report
npm install
\end{lstlisting}

\subsection{Environment Variables}

Create \texttt{.env.local} in the repository root:

\begin{lstlisting}
NEXT_PUBLIC_SUPABASE_URL=https://your-project.supabase.co
NEXT_PUBLIC_SUPABASE_ANON_KEY=your-anon-key-here
OLLAMA_BASE_URL=http://127.0.0.1:11434
\end{lstlisting}

No cloud LLM keys are required to run any of the three features in
this submission.

\subsection{Local Models}

In one terminal:

\begin{lstlisting}
ollama serve
\end{lstlisting}

In another terminal pull the five models the dropdown lists:

\begin{lstlisting}
ollama pull llama3.2
ollama pull phi3:mini
ollama pull qwen2.5:3b
ollama pull tinyllama
ollama pull gemma2:2b
ollama list
\end{lstlisting}

\texttt{ollama list} should show the five tags above.  Total disk
footprint is roughly nine gigabytes.

\subsection{Run the Application}

\begin{lstlisting}
npm run dev
# then open http://localhost:3000/chat/multi
\end{lstlisting}

You should see a chip dropdown with three models pre-selected.  Type a
short prompt (for example, ``Reply with the single word OK.'')
and press \emph{Ask all selected}.  Three cards should render in
parallel, each labelled with its model name, its Ollama family, the
exact tag, a latency badge, and a copy button.  The fastest
successful response will carry a green ``fastest'' marker.

\subsection{Jasmine Unit Tests}

Run each feature's spec independently:

\begin{lstlisting}
npm run test:jasmine:multi1
npm run test:jasmine:multi2
npm run test:jasmine:multi3
\end{lstlisting}

You can also run all three plus the team's other unit tests with
\texttt{npm test}.

\subsection{Cucumber Acceptance Tests}

The Cucumber suites need both the dev server and the local Ollama
daemon running because the steps drive the live UI:

\begin{lstlisting}
# terminal 1
npm run dev

# terminal 2
npm run test:cucumber:multi1
npm run test:cucumber:multi2
npm run test:cucumber:multi3
\end{lstlisting}

The \texttt{@multi1}/\texttt{@multi2}/\texttt{@multi3} tags scope each
run to one feature so a failure in one does not block the others.

\subsection{Puppeteer Tests}

Same pre-requisites:

\begin{lstlisting}
# terminal 1
npm run dev

# terminal 2
npm run test:puppeteer:multi1
npm run test:puppeteer:multi2
npm run test:puppeteer:multi3
\end{lstlisting}

Each script prints \texttt{[PASS]} / \texttt{[FAIL]} per assertion
and exits with a non-zero code on failure.  Screenshots are written
to \texttt{test-screenshots/parthaped-feature*.png}.

\subsection{Process Recap}

\begin{enumerate}[leftmargin=*]
  \item User Story
  \item Jasmine unit tests for the pure helper module
  \item Implement helper until specs pass
  \item Cucumber scenarios against the live UI
  \item Implement React component / API route until scenarios pass
  \item Headed Puppeteer demo as the acceptance gate
  \item Commit and push to \texttt{parthaped-individual-report}
\end{enumerate}

\end{document}
